\section{Related Work}
\label{sec:related}

Our work sits at the intersection of three strands of prior research: RLHF overoptimization theory, bidirectional human-AI alignment measurement, and Goodhart's Law in machine learning. Each strand illuminates part of the phenomenon we study, but none provides the quantitative characterization we offer---a named divergence construct, a regression slope, and a cross-dataset replication.

\subsection{RLHF Overoptimization and Reward Hacking}

Reinforcement learning from human feedback was introduced as a paradigm for aligning language models with human preferences through a reward model trained on preference annotations \cite{Ouyang2022Training,Christiano2017Deep}. RLHF has become the dominant alignment approach, enabling models that score high on helpfulness and harmlessness benchmarks \cite{Bai2022Training,Stiennon2020Learning}. However, the reward model is an imperfect proxy for human preference, and sustained optimization against this proxy causes the policy to exploit proxy weaknesses---a phenomenon known as reward hacking or specification gaming \cite{Krakovna2020Specification,Skalse2022Defining}.

Coste et al.~\cite{Coste2023Reward} provide the clearest empirical demonstration: as KL budget from the base policy increases from 0 to 8 nats, the reward model score rises monotonically while gold human preference peaks and then reverses. Gao et al.~\cite{Gao2023Scaling} characterize the scaling laws of reward model overoptimization at multiple RM sizes, showing consistent patterns across model scales. Both papers describe the divergence qualitatively---RM rises while gold reverses---but stop short of defining the \emph{normalized divergence gap} ($\text{RM\_norm} - \text{gold\_preference}$) as a regression target or reporting cross-dataset slope comparison. Our work takes this step: we fit OLS regression to the normalized gap in both datasets and report $\beta = 0.1433\ \text{nat}^{-1}$ (Coste) and $\beta = 0.1599\ \text{nat}^{-1}$ (Gao), enabling the first cross-dataset characterization of the calibration-alignment divergence curve slope.

Ziegler et al.~\cite{Ziegler2019Fine} and Stiennon et al.~\cite{Stiennon2020Learning} document early cases where RLHF reward models diverge from human preferences in text summarization. These works establish the qualitative phenomenon; our contribution is quantification and cross-dataset effect size comparison.

\subsection{Bidirectional Human-AI Alignment Measurement}

The ICLR 2025 Workshop on Bidirectional Human-AI Alignment synthesized 400 papers in the alignment literature and found a systematic measurement asymmetry: virtually all major alignment benchmarks---TruthfulQA \cite{Lin2022TruthfulQA}, BBQ \cite{Parrish2022BBQ}, HHH-RLHF, WinoBias \cite{Zhao2018Gender}, HELM \cite{Liang2023Holistic}---measure the AI$\to$Human direction exclusively. Human$\to$AI alignment is studied in a parallel HCI and cognitive science track but has not been integrated with ML evaluation frameworks.

Lai et al.~\cite{Lai2021Understanding} demonstrate that AI-assisted decision-making produces over-reliance: in several experimental settings, 60--80\% of humans agree with AI predictions regardless of AI accuracy. Bansal et al.~\cite{Bansal2021Does} and Bu\c{c}inca et al.~\cite{Buccinca2021Trust} further characterize over-reliance patterns and propose interventions. However, these behavioral studies do not connect to RLHF optimization pressure---they measure human calibration at a fixed model deployment snapshot, not as a function of how aggressively the model was optimized.

Our work bridges these two tracks: we treat RLHF overoptimization as the empirical instantiation of bidirectional alignment tension. We scope this as ``evaluation calibration divergence''---distinct from but motivating the study of user behavioral calibration (appropriate reliance; Lai et al.~\cite{Lai2021Understanding} paradigm). The connection to user over-reliance remains a theoretical extension for future behavioral work.

\subsection{Goodhart's Law and Proxy-Target Decoupling in ML}

Goodhart's Law---``when a measure becomes a target, it ceases to be a good measure''---has been formalized in the ML setting by Manheim and Garrabrant~\cite{Manheim2019Categorizing} and others. Krakovna et al.~\cite{Krakovna2020Specification} compile a specification gaming list documenting empirical cases in RL. Skalse et al.~\cite{Skalse2022Defining} formalize reward tampering and misspecification conditions under which proxy-target decoupling is guaranteed.

These theoretical works predict proxy-gold divergence under optimization pressure. Our work provides an empirical instantiation with specific quantitative parameters: in the RLHF setting with explicit reward model training, the proxy-gold gap grows at $\beta \approx 0.143$--$0.160\ \text{nat}^{-1}$ of KL budget. This connects the abstract Goodhart/specification-gaming literature to a concrete, measurable phenomenon in modern RLHF training.

Our three quantitative contributions---the calibration-alignment divergence construct, the regression slope $\beta$ with cross-dataset replication, and the normalized gap metric---are not derivable from any prior paper. The closest works are Coste et al.~\cite{Coste2023Reward} and Gao et al.~\cite{Gao2023Scaling}, which provide the raw experimental data; we contribute the regression framework, normalization methodology, cross-dataset comparison, and bidirectional alignment framing.
